\chapter{Graph algorithms}

\section{Definitions}

Let $V$ be a finite set, and denote by
\[
E(V) = \lbrace \lbrace u, v \rbrace \,|\, u, v \in V, u \ne v \rbrace,
\]
the subsets of $V$ of two distinct elements.
We define a graph on $V$ as below.
\begin{defn}{Graph}
\index{graph}
A pair $G = (V,E)$ with $E \subseteq E(V)$ is called a graph (on $V$).
The elements of $V$ are the vertices, and those of $E$ the edges of the graph.
The \Index{vertex} set of a graph $G$ is denoted by $V_G$ and its edge set by $E_G$.
Therefore $G = (V_G, E_G)$.
\end{defn}

In the literature, graphs are also called simple graphs; vertices are called nodes or points; edges are called lines or links.
A pair $\lbrace u, v \rbrace$ is usually written simply as $uv$.
Notice that then $uv$ = $vu$.
In order to simplify notations, we also write $v \in G$ instead of $v \in V_G$.

\section{Computer representation}

As any graph is basically constituted by nodes and edges, we first define a type for the nodes and a type for the edges.
The graph itself will be represented by an object of type \verb|GR_Graph|.
Various computer representations are possible, but in many applications, graph are sparse,
as many possible edges, corresponding to an ordered pair of nodes (in the case of an oriented
graph) do not exist.
To account of this situation, we decide to represent a graph simply as a collection of nodes, and each node is associated to a set of input edges and a set of output edges.
This representation is well suited for oriented graphs, while for non-oriented graphs, we then need to duplicate each edge.
A graph can then be allocated using the function
\begin{verbatim}
GR_Graph *gr_graph_new(int n)
\end{verbatim}
where $n$ is the number of nodes. to be known in advance.
The user has then to add the nodes and the edges.
In order to help node addition, the function \verb|gr_graph_new| also initialize the set of nodes and give (unique) identifiers from $0$ to $n-1$.
It is possible to force the numbering to start from a different number using the function
\begin{verbatim}
GR_Graph *gr_graph_new_offset(int n, int offset)
\end{verbatim}
where \verb|offset| is the first number to use.
The memory can be freed using the function
\begin{verbatim}
void gr_graph_free(GR_graph *graph)
\end{verbatim}
This function with also free the memory associated to each node of the graph.

The basic type for a node is \verb|GR_Node|, which has the following attributes:
\begin{enumerate}
\item
\verb|id|: a numeric identifier;
\item
\verb|internal_id|: an internal identifier, used for library functions purposes;
\item
\verb|label|: a string label;
\item
\verb|value|: an array of real values attached ti this node;
\item
\verb|in|: the list of input edges;
\item
\verb|out|: the list of output edges.
\end{enumerate}
A node can be created using the function
\begin{verbatim}
GR_Node *gr_node_new()
\end{verbatim}
and freed with
\begin{verbatim}
void gr_node_free(GR_Node *node)
\end{verbatim}

The identifiers can be recovered using the functions
\begin{verbatim}
int gr_node_get_id(GR_Node *node)
\end{verbatim}
and
\begin{verbatim}
int gr_node_get_internal_id(GR_Node *node)
\end{verbatim}
It is also possible the obtain the internal identifier of the node with the identifier \verb|id| in a given graph with the function
\begin{verbatim}
int gr_graph_get_node_internal_id(GR_Graph *graph, int id)
\end{verbatim}
The search is linear, and can be quite expensive in large graphs.
If no node with identifier \verb|id| has been found, the function returns the value -1.

Once the set of nodes created, we can connect them by creating edges.
The basic function creating an edge from node \verb|origin| to node \verb|destination| is
\begin{verbatim}
GR_Edge *gr_edge_new(GR_Node *origin, GR_Node *destination)
\end{verbatim}
The memory associated to an edge can be freed with the function
\begin{verbatim}
void gr_edge_free(GR_Edge *edge)
\end{verbatim}
Each edge has a weight that can be set with the function
\begin{verbatim}
double gr_edge_get_weight(GR_Edge *edge)
\end{verbatim}

An edge can also be created by connecting two nodes of known identifiers in a graph with the function
\begin{verbatim}
GR_Edge *gr_graph_connect(GR_Graph *graph,
                          int origin, int destination,
                          double weight)
\end{verbatim}
As written before, the search of the real identifier of a node can be time consuming, so when possible, it is preferable to use the function below, where \verb|origin| and \verb|destination| are now internal identifiers:
\begin{verbatim}
GR_Edge *gr_graph_create_edge(GR_Graph *graph,
                              int origin, int destination,
                              double weight)
\end{verbatim}
			  
\section{Graphs algorithms}

\subsection{Shortest paths}

Selecting an origin $s$ and a destination $t$ in a graph, the problem is to find a shortest path
between these two nodes.
Various algorithms exist, and the library provides some simple implementations, aiming to
illustrate the techniques.

\subsubsection{Topologically-ordered algorithms}

The first presented algorithms are inspired from dynamic programming techniques.
We assume that the graph is topologically-ordered, i.e. it is ordered and an edge $(i,j)$
can exist only if $i < j$, and that $a_{ij} > 0$ if the edge $(i,j)$ exists.
Three algorithms are proposed.
The first is a direct translation of the dynamic programming algorithm.
If $J(i)$ is the minimum distance from $i$ to $t$, then
\[
 J(j) + a_{ij} = J(i),
\]
if the edge $(i,j)$ belongs to a shortest path.
The function
\begin{verbatim}
void gr_graph_backward_linking(GR_Graph *graph,
                               int origin, int destination,
                               double *J, int *next);
\end{verbatim}
implements the method, fixing the distances from the destination to the origin, represented by their internal identifiers.
On output, \verb|J[i]| contains the distance from node $i$ to the node \verb|destination|,
and \verb|next[i]| contains the successor of node $i$ on the shortest path from
\verb|origin| to \verb|destination|.

\begin{verbatim}
void gr_graph_forward_linking(GR_Graph *graph,
                              int origin, int destination,
                              double *D, int *previous);
\end{verbatim}

\begin{verbatim}
void gr_graph_accession(GR_Graph *graph,
                        int origin, int destination,
                        double *D, int *previous);
\end{verbatim}

\subsubsection{Label correcting algorithm}

\begin{verbatim}
void gr_graph_dijkstra(GR_Graph *graph,
                       int origin, int destination,
                       double *dist, int *previous);
\end{verbatim}

\begin{verbatim}
void gr_graph_dijkstra_edges(GR_Graph *graph,
                             int origin, int destination,
                             double *dist, GR_Edge **previous,
                             double (*weight)(GR_Edge *edge));
\end{verbatim}

The more general algorithm allows to consider edges of negative weight, as long as there is no negative cycle.
\begin{verbatim}
int gr_graph_labelcorrecting_adj(int n,
                                 int origin, int destination,
                                 int **adj, double **v_adj,
                                 double *u,
                                 int *parent, double *w);
\end{verbatim}
